\section{New starting weights test the proposed hybrid}
EXP-007 freezes a follow-up after the full EXP-006 development results. Twelve new starting weights are selected from the same earlier atlas, excluding a closed radius-$0.04$ neighborhood of every development anchor in rectangle-normalized coordinates. The old atlas and its non-nine pixels remain unchanged. The remaining 3,758 eligible pixels are split into margin-rank thirds, with four starts per third; the smallest selected normalized distance from a development origin is $0.04006$. This is a test of new waveform interventions conditional on an observed atlas, not generalization to an unseen map or network. Cross-stratum neighbors and overlap with old intervention destinations remain possible.

All twelve new starts qualify using independent common-initialization integrations at both numerical steps and unchanged primitive-orbit shooting gates. Four new phases use exact indices $337,1012,1687,2362$ of each 2,700-sample endpoint-excluded certified orbit. The corresponding epochs are $iT/2700$; they are approximately, rather than exactly, the odd eighths of a cycle.

\subsection{A stronger comparison with explicit preview access}
The seven policies are atlas margin, signed local sensitivity, signed tangent, signed coordinate axes, uniform 16-direction search, augmented 21-entry search, and a hybrid. The hybrid uses the local method for radii through $0.004$ and atlas margin at $0.016$. This switch is chosen from EXP-006 before any new outcomes. Local sensitivity is integrated from the fixed initial state over $0$--$300$, with its phase-projected Gram restricted to the scored interval below.

Each nonzero condition saves 21 separate simulator previews: 16 uniform directions, both local signs, both tangent signs, and the atlas direction. Duplicates are retained. Each preview starts from the same original orbit state, runs for 75 units at $\Delta t=0.005$, and measures phase-aligned full-activation error. Local and tangent inspect two previews each, axes four, uniform search sixteen, and augmented search all twenty-one. The fixed atlas requires no preview. Augmented search therefore uses geometry and local derivatives; uniform search is the atlas-free search control. All policies use the same native weight-change norm. Their information and compute costs differ.

Policy choices are fixed before independent 300-unit target trajectories at both integration steps. The primary metric scores only the 2,250 samples $75.1\le t\le300$, fitting one constant phase shift within that interval. The preview and scoring samples are disjoint but dynamically dependent. Clonable simulator access is an explicit assumption. Full $0$--$300$ error is secondary. Neither frequency refitting nor time warping is allowed. Equal phase, dose and anchor weighting, missingness rules, and 10,000 descriptive paired-anchor bootstrap draws are frozen in advance. The hybrid and local policy coincide at the three smaller doses, so their pooled contrast is exactly one quarter of the largest-dose contrast; these are not separate demonstrations.

\subsection{The new comparison does not establish a general hybrid advantage}
All 4,032 preview entries and 2,784 target records complete, with no replacement or missing outcome. Table~\ref{tab:holdout-means} reports the primary tail error. At radius $0.016$, hybrid mean error is $0.115448$, versus $0.260336$ for signed local sensitivity. The mean paired reduction is $0.144887$, but its descriptive 95\% interval spans $[-0.045524,0.353194]$ and only six of twelve anchors favor hybrid. The point estimate is 55.65\% lower, not a resolved advantage across these starts. Augmented search has a smaller largest-dose mean, $0.076644$; its paired comparison with hybrid is also unresolved.

At the smallest two radii, local/hybrid beats the tangent and coordinate rules at every anchor. At radius $0.004$, however, its mean error is $0.132630$, versus $0.027130$ for tangent and coordinate policies. Thus the frozen switch fails to capture an intermediate-dose opportunity. This observation does not authorize moving the switch and claiming success on the same data. The pooled hybrid mean is $0.062373$, compared with $0.098595$ for local, $0.051095$ for tangent/axes, $0.074171$ for uniform search and $0.049246$ for augmented search. Every one of the five primary pooled intervals includes zero (Table~\ref{tab:holdout-contrasts}).

\begin{table}[H]
\centering\small
\setlength{\tabcolsep}{4pt}
\begin{tabular}{rrrrrrrr}
\toprule
Radius & Atlas & Local & Tangent & Axes & Uniform & Augmented & Hybrid\\
\midrule
0.00025 & 0.013149 & 0.000244 & 0.002412 & 0.002412 & 0.001440 & 0.000580 & 0.000244\\
0.001 & 0.056555 & 0.001169 & 0.008732 & 0.008732 & 0.005639 & 0.002351 & 0.001169\\
0.004 & 0.135219 & 0.132630 & 0.027130 & 0.027130 & 0.107517 & 0.117410 & 0.132630\\
0.016 & 0.115448 & 0.260336 & 0.166105 & 0.166105 & 0.182088 & 0.076644 & 0.115448\\
All doses & 0.080093 & 0.098595 & 0.051095 & 0.051095 & 0.074171 & 0.049246 & 0.062373\\
\bottomrule
\end{tabular}
\caption{\textbf{EXP-007: all seven frozen policies.} Mean all-neuron activation RMSE on $75<t\le300$, primary $\Delta t=0.005$. Uniform uses sixteen previews; augmented uses twenty-one entries including local and atlas information. All doses are equally weighted. The table retains unfavorable doses and identical-action controls.}
\label{tab:holdout-means}
\end{table}

\begin{table}[H]
\centering\small
\setlength{\tabcolsep}{4pt}
\begin{tabular}{lrr}
\toprule
Comparator & Largest-dose benefit [95\% interval] & Pooled benefit [95\% interval]\\
\midrule
Local & $0.144887\;[-0.045524,0.353194]$ & $0.036222\;[-0.011381,0.088298]$\\
Tangent & $0.050656\;[-0.043143,0.153512]$ & $-0.011278\;[-0.050720,0.028091]$\\
Axes & $0.050656\;[-0.043143,0.153512]$ & $-0.011278\;[-0.050720,0.028091]$\\
Uniform 16 & $0.066640\;[-0.042759,0.201205]$ & $0.011798\;[-0.028640,0.052582]$\\
Augmented 21 & $-0.038805\;[-0.097362,0.031163]$ & $-0.013127\;[-0.033383,0.008182]$\\
\bottomrule
\end{tabular}
\caption{\textbf{All declared comparator-minus-hybrid contrasts.} Positive values favor hybrid. Intervals resample twelve selected anchor clusters and are descriptive; phases and integration steps are not independent replications. All ten displayed intervals include zero. No resolved overall advantage or equal-total-cost benefit is established.}
\label{tab:holdout-contrasts}
\end{table}

\subsection{A coordinate control explains the tangent result in this sample}
The tangent and signed-axis policies choose exactly the same actions in all 192 nonzero conditions: $+\gamma$ in 168 and $-\gamma$ in 24, always holding $\alpha$ fixed. All 192 corresponding preview state arrays and all 384 target state arrays across both steps are identical. The old nearest-pixel normals happen to be horizontal at all twelve selected starts, so their tangents are coordinate directions. This equivalence is a result of the retained action records, not a renaming of a comparator after outcomes. Tangent performance here therefore supplies no independent evidence that detailed shrimp geometry is necessary; a simple signed coordinate test realizes the same actions. It does not establish that changing $\gamma$ is always best or that geometry has no value elsewhere.

The preselected example also prevents an exclusively favorable illustration. At the first eligible start and phase, largest-dose hybrid error is $0.059668$, local error $0.037763$, and both uniform and augmented search error $0.589753$. The augmented selector minimizes its preview objective but does not minimize later waveform error. The example is retained in Figure~\ref{fig:holdout-waveforms}; its ranking differs from the aggregate largest-dose ranking. More previews are not automatically a better predictor when the selection and evaluation horizons differ.

\subsection{Numerical and cost limits}
All 48 fine-step zero-dose shams complete, with mean tail error $6.53\times10^{-5}$ and maximum $6.88\times10^{-5}$. Across 1,392 paired coarse/fine targets, the absolute RMS difference has median $1.30\times10^{-7}$, mean $0.00403$, 95th percentile $0.01456$ and maximum $0.23486$. The coarse results preserve the unresolved pooled-comparison conclusion, but these differences rule out a claim of uniform pointwise convergence.

Collection and summary construction take 137.51 seconds with four local workers, excluding final verification. Forty-eight sensitivity solves consume 23.52 seconds within that total. Summing the overlapping preview-call wall times gives 161.85 seconds; that is neither elapsed duration nor CPU time. The old atlas costs a separate 140.54 seconds, and per-policy preview counts range from zero to twenty-one. Shared research execution does not make these information sources free. All 7,149 receipt-listed artifacts are retained. Independent interpolation reconstructs 504,320 saved RMS values across 9,600 metric blocks with maximum absolute discrepancy $4.44\times10^{-16}$. No paid resource is used.


\clearpage
\newgeometry{margin=0.55in}
\begin{landscape}
\begin{figure}[H]
\centering
\includegraphics[width=\linewidth,height=0.78\textwidth,keepaspectratio]{neuralshrimp-holdout-anchors.pdf}
\caption{\textbf{Where the new interventions start.} The measured EXP-004 period atlas underlies the new and development anchor locations. The second view expresses coordinates in units of the focus rectangle widths so the fixed radius-$0.04$ exclusion neighborhoods can be shown as circles. The exclusion applies to starting locations, not every previous destination or trajectory. All twelve attempted new starts are retained. The same previously observed atlas supplies margins and tangents; the design does not test a new network or an unseen parameter map.}
\label{fig:holdout-anchors}
\end{figure}
\end{landscape}
\restoregeometry

\clearpage
\newgeometry{margin=0.55in}
\begin{landscape}
\begin{figure}[H]
\centering
\includegraphics[width=\linewidth,height=0.78\textwidth,keepaspectratio]{neuralshrimp-holdout-utility.pdf}
\caption{\textbf{New-start policy comparison with stronger controls.} Dose means at both integration steps score only $75<t\le300$. Fine-step preview choices are reused at both steps. Positive comparator-minus-hybrid contrasts favor the hybrid; all five declared comparator contrasts are retained, including the largest-dose per-anchor results. Intervals are descriptive paired-anchor bootstrap percentiles, not population confidence claims. The preview budget panel counts candidate calls, not equal total resources: atlas construction and local sensitivity are additional dependencies. Augmented 21-entry search includes atlas and local directions; uniform 16-direction search does not. The hybrid copies local at three doses, making those local contrasts identically zero.}
\label{fig:holdout-utility}
\end{figure}
\end{landscape}
\restoregeometry

\clearpage
\newgeometry{margin=0.55in}
\begin{landscape}
\begin{figure}[H]
\centering
\includegraphics[width=\linewidth,height=0.78\textwidth,keepaspectratio]{neuralshrimp-holdout-waveforms.pdf}
\caption{\textbf{A prospectively selected new-start waveform example.} First eligible anchor, first phase, largest dose, and primary numerical step. Rows show actual activations and componentwise errors. Preview time is shaded; the primary endpoint excludes that interval. Candidate curves are shown against the original synchronous reference, while each error curve uses its own single phase shift fitted to the entire scored tail. No frequency correction or separate cycle alignment is used. This example follows the frozen selection rule even if it is unfavorable to the hybrid or unrepresentative of the aggregate.}
\label{fig:holdout-waveforms}
\end{figure}
\end{landscape}
\restoregeometry
